\chapter{Inequalities between Matrices}\label{linear:ch:inequalities}

\begin{center}\itshape
All analysts spend half their time hunting through the literature for inequalities which they want to use and cannot prove.
\upshape(G.~H.~Hardy, \emph{Prolegomena to a chapter on inequalities}, 1929)
\end{center}
\noindent\emph{Use the spectral results of Chapter~\ref{linear:ch:spectral} and inner-product tools of Appendix~\ref{linear:ch:inner}. This appendix is optional for the main chapters; Volumes~12 and~15 use its Hadamard bound. Matrix sizes below are positive.}

Hadamard asked how large the determinant of a real matrix can be when every entry lies in $[-1,1]$. Each column has length at most $\sqrt n$, and their parallelepiped has greatest volume when the columns are perpendicular. This gives $|\det A|\le n^{n/2}$. Equality requires entries $\pm1$ and orthogonal columns, a Hadamard matrix.

His 1893 work included constructions at orders twelve and twenty. For orders at least three, divisibility by four is necessary. Existence at every positive multiple of four remains the general Hadamard conjecture. There has been recent progress worth dating precisely: Alpöge, Voinov and Reynolds-Haertle announced constructions in August 2026 at the twelve formerly open orders below 2000, including 668. Egan, Kotsireas, London and \'O Cath\'ain verified them in arXiv:2610.27175, first submitted on 19 October and revised on 3 November 2026. Their report establishes existence at every positive multiple of four below 2000. Volume~3, Appendix~D, gives the earlier Sylvester and Paley constructions.

Comparing Hermitian matrices means comparing their quadratic forms. Congruence preserves that order, inversion reverses it, and square roots preserve it. Squaring usually fails. We will use these distinctions for determinants, block matrices, covariance and kernels, then finish with eigenvalue discs and Hilbert's inequality.

\begin{think}
If $A\succeq B\succeq0$, can squaring reverse some quadratic-form comparison? Try two positive ellipses with different principal directions. Commuting matrices reduce to scalar inequalities; a counterexample must use their failure to commute.
\end{think}

\section{The Loewner order}

A Hermitian matrix is positive semidefinite, $A\succeq0$, when every quadratic form $x^*Ax$ is non-negative, equivalently all eigenvalues are non-negative. Strict positivity for nonzero $x$ means positive definiteness, $A\succ0$. Define $A\succeq B$ by positivity of $A-B$. Quadratic forms show reflexivity, transitivity and antisymmetry, making this a partial order. Loewner studied its function theory in 1934.

\begin{proposition}\label{linear:prop:loewner}
For Hermitian matrices of the same size:
\begin{enumerate}[label=(\roman*)]
\item $A\succeq B$ implies $X^*AX\succeq X^*BX$ for every rectangular $X$ of compatible size.
\item It also implies $\lambda_k(A)\ge\lambda_k(B)$ in matching eigenvalue order, hence the trace inequality and, if $B\succeq0$, $\det A\ge\det B$.
\item $A\succeq B\succ0$ implies $B^{-1}\succeq A^{-1}\succ0$.
\item For positive definite $A,B$, the order $A\succeq B$ is equivalent to $E_A\subseteq E_B$, where $E_A=\{x:x^*Ax\le1\}$.
\end{enumerate}
\end{proposition}
\begin{proof}
(i) Test on $y$: the quadratic form is $(Xy)^*(A-B)(Xy)\ge0$. (ii) Courant--Fischer (Theorem~\ref{linear:thm:courantfischer}) compares each min--max Rayleigh quotient. Summing or, for non-negative eigenvalues, multiplying gives the consequences.

(iii) The positive square root and inverse are already available from main Chapters~\ref{linear:ch:spectral} and~\ref{linear:ch:svd}. Set $C=B^{-1/2}AB^{-1/2}\succeq I$. Its eigenvalues are at least one, so $C^{-1}\preceq I$. Congruence by $B^{-1/2}$ gives $A^{-1}\preceq B^{-1}$. (iv) The forward direction is the quadratic-form inequality. If the order fails, choose $x$ with $x^*Bx>x^*Ax>0$ and scale it so the first form exceeds one while the second is below one. This contradicts containment.
\end{proof}

Larger positive forms therefore give smaller ellipsoids. For an exact squaring counterexample, take
\[A=\begin{pmatrix}2&1\\1&1\end{pmatrix},\qquad B=\operatorname{diag}(9/10,1/20).
\]
The difference $A-B$ has positive first diagonal entry and determinant $9/200$, so it is positive definite by completing the square. But $A^2-B^2$ has determinant $-25219/40000<0$. Its two real eigenvalues have opposite signs, so it is not positive semidefinite. Figure~\ref{linear:fig:loewner} shows the ellipsoid comparison.
\begin{figure}[ht]
\centering
\fig{linear-ch11-loewner}
\caption{Left: $A\succeq B$ means that the ellipse $E_A$ (charcoal) lies inside $E_B$ (rust, here a long ellipse, clipped). Right:
$E_{A^2}$ is thinner and longer along a tilted axis, and leaves $E_{B^2}$: $A^2\not\succeq B^2$. Squaring is not monotone for matrices.}
\label{linear:fig:loewner}
\end{figure}

\begin{theorem}[L\"owner--Heinz for the square root]\label{linear:thm:sqrtmonotone}
For $A\succeq B\succeq0$, the positive square roots satisfy $\sqrt A\succeq\sqrt B$.
\end{theorem}
\begin{proof}
Put $S=\sqrt A,T=\sqrt B$. If $(S-T)x=\lambda x$ for a unit eigenvector, Hermitian symmetry gives
\[0\le x^*(A-B)x=x^*(S(S-T)+(S-T)T)x
 =\lambda(x^*Sx+x^*Tx).
\]
If $\lambda<0$, both non-negative quadratic terms must vanish. For a positive matrix, $x^*Sx=0$ implies $Sx=0$ by its spectral decomposition, and likewise $Tx=0$. This contradicts $(S-T)x=\lambda x$. Thus every eigenvalue of the difference is non-negative.
\end{proof}

\begin{remark}
On a real interval $J$, call $f$ operator monotone if applying spectral calculus preserves this order in every dimension for spectra in $J$. Heinz proved monotonicity of $t^p$ on $[0,\infty[$ for $0\le p\le1$. On an open interval, Loewner's theorem characterises operator monotone functions by holomorphic extensions to the upper half-plane with non-negative imaginary part, extending analytically across $J$ with real values there. Constant real functions are included. The logarithm belongs to the theory on $(0,\infty)$, not at zero; $t/(1+t)$ also preserves positive-matrix order, while $t^2$ and $e^t$ do not in all dimensions. We quote this characterisation; only the square-root case was proved above.
\end{remark}

\section{Schur complements}

For Hermitian $M=\begin{pmatrix}A&B\\B^*&C\end{pmatrix}$ with invertible $A$, define its Schur complement $M/A=C-B^*A^{-1}B$.
\begin{proposition}[Schur complements]\label{linear:prop:schurcomplement}
With $X=\begin{pmatrix}I&-A^{-1}B\\0&I\end{pmatrix}$, one has
\[X^*MX=\operatorname{diag}(A,C-B^*A^{-1}B).
\]
Consequently $\det M=\det A\det(M/A)$. If $A\succ0$, then $M\succeq0$ exactly when $M/A\succeq0$, and strict positivity is likewise equivalent.
\end{proposition}
\begin{proof}
Multiplication first kills the upper-right block, and multiplication by $X^*$ then kills the lower-left block. Since $X$ is invertible with determinant one, determinants agree and congruence in both directions preserves positivity. A block diagonal form is positive precisely when both diagonal blocks are, proving both tests.
\end{proof}

\begin{proposition}[Sylvester's determinant identity]\label{linear:prop:detIAB}
For complex $A\in M_{m,n},B\in M_{n,m}$,
\[\det(I_m+AB)=\det(I_n+BA).
\]
In particular $\det(I+uv^*)=1+v^*u$.
\end{proposition}
\begin{proof}
Compute the determinant of $\begin{pmatrix}I_m&-A\\B&I_n\end{pmatrix}$. Eliminating its lower-left block using the first block row gives diagonal blocks $I_m,I_n+BA$. Eliminating by a block column operation instead gives $I_m+AB,I_n$. Both elimination matrices have determinant one, so the two values agree. Taking the inner dimension to be one gives the last formula.
\end{proof}

This replaces a large determinant by a small one when the inner dimension is small. Applying the same polynomial identity with a scalar parameter also recovers matching nonzero eigenvalues of $AB$ and $BA$. Exercise~\ref{linear:ex:sherman} uses the rank-one form to update an inverse.

\section{Determinant inequalities}
\begin{theorem}[Hadamard's inequality]\label{linear:thm:hadamard}
A positive semidefinite matrix satisfies $\det A\le\prod_i a_{ii}$. For positive definite $A$, equality requires $A$ diagonal.
\end{theorem}
\begin{proof}
A zero diagonal entry forces its column to vanish by spectral positivity, making both sides zero. Otherwise normalise by $D=\operatorname{diag}(a_{ii}^{-1/2})$. The positive matrix $C=DAD$ has unit diagonal and eigenvalue sum $n$. AM--GM gives $\det C\le1$, proving the bound. Equality forces all eigenvalues to equal one, so $C=I$ and $A=D^{-2}$ is diagonal.
\end{proof}

\begin{corollary}[Hadamard's volume inequality]\label{linear:cor:hadamardvol}
For square $X$ with columns $x_i$, $|\det X|\le\prod_i\|x_i\|$. Equality holds exactly when the columns are pairwise orthogonal or at least one is zero. In particular real entries in $[-1,1]$ give $|\det X|\le n^{n/2}$.
\end{corollary}
\begin{proof}
Apply Hadamard to $X^*X$, whose determinant is $|\det X|^2$. With no zero columns, a singular $X$ gives strict inequality. An invertible equality case makes the Gram matrix diagonal, hence the columns orthogonal. Both listed cases do attain equality. The entry bound makes each column norm at most $\sqrt n$.
\end{proof}

Hadamard matrices satisfy $HH^{\mathsf T}=nI$ with entries $\pm1$. The Walsh character tables on $\mathbb F_2^k$ give examples of order $2^k$; see Chapter~\ref{ch:finite} and the Sylvester construction in Volume~3. They concern the binary additive group, rather than ordinary cyclic convolution. Exercise~\ref{linear:ex:hadamard4} proves the divisibility restriction.

\begin{theorem}[Fischer's inequality]\label{linear:thm:fischer}
For a positive semidefinite block matrix $M=\begin{pmatrix}A&B\\B^*&C\end{pmatrix}$ with square diagonal blocks, $\det M\le\det A\det C$.
\end{theorem}
\begin{proof}
If $A$ is singular, a nonzero $x\in\ker A$ makes $(x,0)$ have zero quadratic form under $M$, hence belong to its kernel. Both sides are then zero. Otherwise $A\succ0$ and $0\preceq M/A\preceq C$. Determinant monotonicity and the Schur determinant formula give the result.
\end{proof}

\begin{theorem}[Minkowski's determinant inequality]\label{linear:thm:minkowskidet}
For positive semidefinite matrices of size $n\ge1$,
\[\det(A+B)^{1/n}\ge\det(A)^{1/n}+\det(B)^{1/n}.\]
\end{theorem}
\begin{proof}
If both are singular, the right side is zero. Otherwise, interchange them if needed so $A\succ0$. Congruence scales all three determinant roots by the same positive factor. Using $A^{-1/2}$ and a unitary diagonalisation reduces to $A=I,B=\operatorname{diag}(d_i)$ with $d_i\ge0$. AM--GM gives
\[\prod_i(1+d_i)^{-1/n}+\prod_i\left(\frac{d_i}{1+d_i}\right)^{1/n}\le1.
\]
Multiplication by $\prod_i(1+d_i)^{1/n}$ proves the claim.
\end{proof}

Homogeneity turns this superadditivity into concavity of $\det^{1/n}$ on the positive cone. The logarithmic version is often useful too.
\begin{theorem}[concavity of $\log\det$]\label{linear:thm:logdet}
For $A,B\succ0$ and $0\le t\le1$,
\[\log\det(tA+(1-t)B)\ge t\log\det A+(1-t)\log\det B.\]
\end{theorem}
\begin{proof}
Congruence adds the same scalar $\log|\det X|^2$ to each side. Reduce to $B=I,A=\operatorname{diag}(d_i)$ with $d_i>0$. The claim follows by summing scalar concavity $\log(td_i+1-t)\ge t\log d_i$.
\end{proof}

\begin{remark}
A nonsingular real Gaussian covariance $\Sigma$ gives differential entropy $\tfrac12\log\det(2\pi e\Sigma)$. Thus interpolating covariance matrices gives a Gaussian entropy at least the corresponding weighted average of Gaussian entropies. This statement concerns Gaussian laws with the interpolated covariance; a mixture of Gaussian laws need not itself be Gaussian. With variances fixed, Hadamard maximises Gaussian entropy at a diagonal covariance, and Fischer gives the corresponding block inequality. Cover--Thomas develop these information-theoretic interpretations.
\end{remark}
\section{The Schur product theorem}

Ordinary products of positive matrices may fail even to be Hermitian. Entrywise multiplication has a different positivity rule.
\begin{definition}\label{linear:def:hadamardproduct}
For two matrices of the same rectangular size, define their \emph{Hadamard} or \emph{Schur product} by $(A\circ B)_{ij}=a_{ij}b_{ij}$.
\end{definition}

\begin{theorem}[Schur, 1911]\label{linear:thm:schurproduct}
Positive semidefinite $A,B$ have positive semidefinite $A\circ B$. If both are positive definite, so is their Schur product.
\end{theorem}
\begin{proof}
Write $B=\sum_kb_kb_k^*$ by spectral decomposition. Each $A\circ b_kb_k^*=D_kAD_k^*$, with $D_k=\operatorname{diag}(b_k)$, is positive semidefinite. Summation proves the first assertion. If $A,B\succ0$ and $x_i\ne0$ for some coordinate of a nonzero $x$, then $b_{ii}=\sum_k|b_{k,i}|^2>0$. Some $D_k^*x$ is therefore nonzero, making its positive $A$ quadratic form strictly positive. The sum is strictly positive too.
\end{proof}

\begin{corollary}\label{linear:cor:entrywiseexp}
For $A\succeq0$, the matrix $(e^{a_{ij}})$ is positive semidefinite. More generally, entrywise application of a power series $f(z)=\sum_{k\ge0}c_kz^k$, with non-negative coefficients and convergence at every entry, preserves positive semidefiniteness. Consequently $(e^{-\gamma\|x_i-x_j\|^2})$ is positive semidefinite for arbitrary real points and $\gamma>0$.
\end{corollary}
\begin{proof}
Schur powers are positive by induction, starting with the all-ones matrix at degree zero. Non-negative finite sums remain positive, and entrywise limits preserve every finite quadratic-form inequality. For the Gaussian matrix, factor its entry as $e^{-\gamma\|x_i\|^2}e^{2\gamma x_i\cdot x_j}e^{-\gamma\|x_j\|^2}$. The middle matrix is an entrywise exponential of a Gram matrix; the outer factors form a diagonal congruence.
\end{proof}

Positive kernels are precisely functions whose finite evaluation matrices are positive semidefinite. Main Chapter~\ref{ch:hilbert} constructs their feature Hilbert spaces. Positive sums, products and pointwise limits give new kernels by the preceding arguments. For distinct points, the Gaussian matrix is actually positive definite; Exercise~\ref{ex:v8cEb} proves this stronger assertion using Fourier inversion and a Vandermonde argument.

\section{Gershgorin's discs}
\begin{theorem}[Gershgorin, 1931]\label{linear:thm:gershgorin}
Every eigenvalue of complex square $A$ belongs to at least one disc
\[D_i=\{z:|z-a_{ii}|\le r_i\},\qquad r_i=\sum_{j\ne i}|a_{ij}|.\]
\end{theorem}
\begin{proof}
In $Ax=\lambda x$, choose a largest-modulus coordinate $|x_i|>0$. Its row equation gives $|\lambda-a_{ii}||x_i|\le\sum_{j\ne i}|a_{ij}||x_j|\le r_i|x_i|$. Cancel the positive coordinate magnitude.
\end{proof}
\begin{figure}[ht]
\centering
\fig{linear-ch11-gershgorin}
\caption{The four Gershgorin discs of a $4\times4$ complex matrix (centres $a_{ii}$, radii the row sums $r_i$) and its four eigenvalues
(black dots). The isolated disc on the right contains exactly one eigenvalue; the three overlapping discs share the other three.}
\label{linear:fig:gershgorin}
\end{figure}

\begin{corollary}\label{linear:cor:diagdom}
Strict diagonal dominance, $|a_{ii}|>r_i$ for every row, makes a matrix invertible. If it is Hermitian and $a_{ii}>r_i$ for every row, it is positive definite.
\end{corollary}
\begin{proof}
The first condition keeps zero outside every eigenvalue disc. Under the second, every disc lies in the open right half-plane and the eigenvalues are real, hence positive.
\end{proof}

Transposition supplies column-sum discs as well. A positive diagonal similarity redistributes off-diagonal row sizes without changing eigenvalues, often improving a bound as in Exercise~\ref{linear:ex:gershscale}. The stronger component-counting form says a union of $k$ discs separated from the others contains $k$ eigenvalues with multiplicity. We quote that extension, illustrated by Figure~\ref{linear:fig:gershgorin}; its proof uses continuity of polynomial roots, while the basic containment just proved suffices below.

\section{A gem: Hilbert's inequality and the number $\pi$}

Finite Hilbert matrices are positive Gram matrices with very small least eigenvalues. Their largest eigenvalues behave differently: a common upper bound survives in every size.

\begin{theorem}[Hilbert's inequality]\label{linear:thm:hilbertineq}
For real square-summable sequences indexed from zero, the double series converges absolutely and satisfies
\[\left|\sum_{m,n\ge0}\frac{a_mb_n}{m+n+1}\right|
 \le\pi\left(\sum_ma_m^2\right)^{1/2}\left(\sum_nb_n^2\right)^{1/2}.
\]
Equivalently, every finite Hilbert matrix has $\|H_N\|_2=\lambda_{\max}(H_N)\le\pi$.
\end{theorem}
\begin{proof}
Set $p_n=(n+1/2)^{-1/2}$ and $k_{mn}=1/(m+n+1)$. First we prove $\sum_nk_{mn}p_n\le\pi p_m$. The function $f(x)=((m+1+x)\sqrt{x+1/2})^{-1}$ is convex on $x>-1/2$: each factor is positive, decreasing and convex, and the product's second derivative has three non-negative terms. Its midpoint value is at most its interval average. At the first interval's singular endpoint use symmetry about the midpoint and an improper limit; the integral is finite. Summing gives
\[\sum_{n\ge0}k_{mn}p_n\le\int_0^\infty\frac{\dd y}{(m+1/2+y)\sqrt y}
 =\frac1{\sqrt{m+1/2}}\int_0^\infty\frac{2\dd u}{1+u^2}=\pi p_m.
\]
Now apply Cauchy--Schwarz on a finite rectangle to the factors $\sqrt{k_{mn}p_n/p_m}|a_m|$ and $\sqrt{k_{mn}p_m/p_n}|b_n|$. The weighted row bound and symmetry give
\[\sum_{m,n}k_{mn}|a_m||b_n|\le\sqrt{\pi\sum_ma_m^2}\sqrt{\pi\sum_nb_n^2}.
\]
Increasing finite rectangles proves absolute convergence and the inequality. For finite sections this bounds the bilinear operator norm; positivity makes it the largest eigenvalue. Conversely uniform finite-section bounds imply the sequence version by the same truncation argument.
\end{proof}

Exercise~\ref{ex:v8cEd} proves that $\pi$ is optimal, so the largest eigenvalues increase to it. Schur established sharpness in 1911. De Bruijn--Wilf's 1962 asymptotic, quoted here, gives the gap $\pi-\lambda_{\max}(H_N)\sim\pi^5/(2\log^2N)$. This is slow convergence; numerical eigenvalues at moderate sizes illustrate it but do not prove either sharpness or the asymptotic. Hilbert's lectures, published by Weyl in 1908, had given constant $2\pi$. The integral counterpart with kernel $1/(x+y)$ is developed in Hardy--Littlewood--P\'olya.

\begin{beauty}
Congruence lets us normalise one positive matrix and diagonalise the other. Scalar AM--GM and concavity then control determinants of matrices that did not commute originally. Schur products use diagonal congruences instead. Hilbert's infinite array needs a different comparison: the weights $(n+1/2)^{-1/2}$ turn row sums into the integral $\int_0^\infty2/(1+u^2)\dd u=\pi$.
\end{beauty}

\begin{figure}[ht]
\centering\fig{appE-rw-projection}
\caption{A direction can be chosen so that finitely many distinct points have distinct scalar projections. Restricting their Fourier exponential sum to that direction gives a Vandermonde system, proving strict Gaussian-kernel positivity.}
\label{fig:v8appEprojection}
\end{figure}
\begin{figure}[ht]
\centering\fig{appE-rw-hilbert}
\caption{In logarithmic coordinates $v=\log x,w=\log y$, the shaded region has $v\in[\log M,\log L-\log M]$ and $w-v\in[-\log M,\log M]$. It stays inside the square and provides a lower bound approaching Hilbert's constant.}
\label{fig:v8appEhilbert}
\end{figure}
\section{Exercises}
\begin{exercise}[$\star$]\label{linear:ex:commutingsquare}
Prove $A^2\succeq B^2$ when $A\succeq B\succeq0$ and $AB=BA$.
\end{exercise}
\begin{exercise}[$\star$]\label{linear:ex:blockcs}
For a positive block matrix $\begin{pmatrix}A&B\\B^*&C\end{pmatrix}$, establish $|x^*By|^2\le(x^*Ax)(y^*Cy)$ for all compatible vectors.
\end{exercise}
\begin{exercise}[$\star\star$]\label{linear:ex:sherman}
For invertible $A$ and $c=1+v^*A^{-1}u\ne0$, prove $\det(A+uv^*)=c\det A$ and $(A+uv^*)^{-1}=A^{-1}-A^{-1}uv^*A^{-1}/c$. Estimate the arithmetic cost when the old inverse is known.
\end{exercise}
\begin{exercise}[$\star\star$]\label{linear:ex:hadamard4}
For a real order-$n$ Hadamard matrix with $n\ge3$, normalise the first row to all ones using column sign changes. Count the four sign patterns of rows two and three to prove $4\mid n$.
\end{exercise}
\begin{exercise}[$\star\star$]\label{linear:ex:maxdet3}
Prove that any sign matrix has determinant divisible by $2^{n-1}$. Combine this with Hadamard's bound to find the maximum absolute determinant at size three, and exhibit a matrix attaining four.
\end{exercise}
\begin{exercise}[$\star$]\label{linear:ex:minkowskieq}
For $A,B\succ0$, classify equality in Minkowski's determinant inequality: prove it occurs exactly when $B=cA$ with $c>0$.
\end{exercise}
\begin{exercise}[$\star\star$]\label{linear:ex:logdetderiv}
For $A\succ0$ and Hermitian $H$, prove the first derivative at zero of $\log\det(A+tH)$ is $\operatorname{tr}(A^{-1}H)$ and its second is $-\operatorname{tr}((A^{-1}H)^2)\le0$. Recover concavity along segments.
\end{exercise}
\begin{exercise}[$\star\star$]\label{linear:ex:schurmin}
Show $\lambda_{\min}(A\circ B)\ge\lambda_{\min}(A)\min_ib_{ii}$ for positive semidefinite $A,B$.
\end{exercise}
\begin{exercise}[$\star\star$]\label{linear:ex:kernels}
Prove positivity of the four evaluation matrices with entries $\min(s_i,s_j)$ for $s_i\ge0$, $(x_i+x_j)^{-1}$ for $x_i>0$, $\cos(\theta_i-\theta_j)$ for real angles, and $(1-u_i\cdot u_j)^{-1}$ for real vectors of norm below one.
\end{exercise}
\begin{exercise}[$\star$]\label{linear:ex:gershscale}
Let $A$ have diagonal $1,2,3$ and all off-diagonal entries $\varepsilon$, with $0<\varepsilon<1/4$. Weyl bounds its smallest eigenvalue $\mu$ within $2\varepsilon$ of one. Use a positive diagonal similarity and Gershgorin to improve this to $|\mu-1|\le8\varepsilon^2$.
\end{exercise}
\begin{exercise}[$\star\star$]\label{linear:ex:diagdomdet}
Prove that a strictly diagonally dominant real matrix with positive diagonal has positive determinant, by a deformation keeping it invertible.
\end{exercise}
\begin{exercise}[$\star\star\star$]\label{linear:ex:oppenheim}
For positive semidefinite $A,B$, prove Oppenheim's inequality $\det(A\circ B)\ge\det A\prod_ib_{ii}$ by induction with Schur complements and Schur products.
\end{exercise}

\begin{exercise}[$\star\star$]\label{ex:v8cEa}
\emph{Strict logarithmic concavity.} Strengthen the log-determinant theorem: for $0<t<1$ and distinct positive definite $A,B$, prove strict inequality. Show that $g_A(H,K)=\operatorname{tr}(A^{-1}HA^{-1}K)$ is a real inner product on Hermitian directions, and relate $g_A(H,H)$ to the second derivative.
\end{exercise}
\begin{exercise}[$\star\star$]\label{ex:v8cEb}
\emph{Distinct Gaussian points.} Prove strict positive definiteness of $(e^{-\gamma\|x_i-x_j\|^2})$ when the real points $x_i$ are distinct and $\gamma>0$. Express its quadratic form using the positive Gaussian Fourier density, then restrict a vanishing exponential sum to a line with distinct projected frequencies.
\end{exercise}
\begin{exercise}[$\star\star$]\label{ex:v8cEc}
\emph{A rank-$r$ inverse update.} For invertible $A\in M_n(\C)$ and $U,V\in M_{n,r}(\C)$, set $C=I_r+V^*A^{-1}U$. Prove $\det(A+UV^*)=\det A\det C$. When $C$ is invertible, establish
\[(A+UV^*)^{-1}=A^{-1}-A^{-1}UC^{-1}V^*A^{-1}.\]
\end{exercise}
\begin{exercise}[$\star\star\star$]\label{ex:v8cEd}
\emph{Why the Hilbert constant is exactly $\pi$.} For $0\le m<N$ put $a_m=(m+1)^{-1/2}$ and zero elsewhere. Prove
\[\sum_{m,n<N}\frac{a_ma_n}{m+n+1}\ge\int_1^{N+1}\int_1^{N+1}\frac{\dd x\dd y}{\sqrt{xy}(x+y)}.\]
For fixed $M>1$ and $N+1>M^2$, restrict to $M\le x\le(N+1)/M$ and $1/M\le y/x\le M$. Divide by $\sum_{m<N}a_m^2$, let $N\to\infty$ and then $M\to\infty$, and deduce $\lambda_{\max}(H_N)\to\pi$.
\end{exercise}

\section{Hints}
\begin{hint}{linear:ex:commutingsquare}
Diagonalise the commuting Hermitian restrictions in a common orthonormal basis.
\end{hint}
\begin{hint}{linear:ex:blockcs}
Compress to the two vectors $(x,0),(0,y)$ and take the two-by-two determinant.
\end{hint}
\begin{hint}{linear:ex:sherman}
Factor out $A$ for the determinant; multiply the proposed inverse by the updated matrix for verification.
\end{hint}
\begin{hint}{linear:ex:hadamard4}
The three row-orthogonality equations force all four sign-pattern counts to agree.
\end{hint}
\begin{hint}{linear:ex:maxdet3}
Adding the first row to every later row exposes $n-1$ factors of two.
\end{hint}
\begin{hint}{linear:ex:minkowskieq}
Equality in the normalised proof requires equal AM--GM inputs.
\end{hint}
\begin{hint}{linear:ex:logdetderiv}
Use the eigenvalues of the Hermitian $A^{-1/2}HA^{-1/2}$.
\end{hint}
\begin{hint}{linear:ex:schurmin}
Schur-multiply $A-\lambda_{\min}(A)I$ by $B$.
\end{hint}
\begin{hint}{linear:ex:kernels}
Use interval indicators, decaying exponentials, planar unit vectors and the geometric power series, respectively.
\end{hint}
\begin{hint}{linear:ex:gershscale}
Choose $D=\operatorname{diag}(1,4\varepsilon,4\varepsilon)$, keeping the two remote discs above the smallest eigenvalue.
\end{hint}
\begin{hint}{linear:ex:diagdomdet}
Scale only the off-diagonal part from zero to its full value.
\end{hint}
\begin{hint}{linear:ex:oppenheim}
The first Schur complement of $A\circ B$ dominates the Schur product of the complement of $A$ with the trailing block of $B$.
\end{hint}
\begin{hint}{ex:v8cEa}
Congruence converts $g_A$ to the trace product on Hermitian matrices.
\end{hint}
\begin{hint}{ex:v8cEb}
A continuous finite exponential sum with zero integral against a strictly positive density vanishes everywhere. Differentiate along a direction with distinct point projections.
\end{hint}
\begin{hint}{ex:v8cEc}
Apply Sylvester's identity at inner size $r$. For the inverse use $(I+V^*A^{-1}U)C^{-1}=I$ in the product expansion.
\end{hint}
\begin{hint}{ex:v8cEd}
On each unit square, the integrand is bounded by the corresponding sum term. Substitute $y=xt$ on the restricted domain; the two integrals separate.
\end{hint}

\section{Solutions}
\begin{solution}{linear:ex:commutingsquare}
\idea Common eigenvectors reduce the order comparison to non-negative scalars.
\writeup A Hermitian $B$ preserves each eigenspace of commuting Hermitian $A$. Diagonalise its Hermitian restriction on each one to obtain a common orthonormal eigenbasis. The order then reads $\alpha_i\ge\beta_i\ge0$, hence $\alpha_i^2\ge\beta_i^2$. Conjugating back proves the claim.
\end{solution}
\begin{solution}{linear:ex:blockcs}
\idea A two-vector compression retains positivity and yields the desired scalar determinant.
\writeup Compress to columns $(x,0),(0,y)$. Its matrix is $\begin{pmatrix}x^*Ax&x^*By\\y^*B^*x&y^*Cy\end{pmatrix}\succeq0$. Its determinant is non-negative, exactly the claimed inequality, including zero-vector cases.
\end{solution}
\begin{solution}{linear:ex:sherman}
\idea The determinant reduces to inner size one, and an explicit rank-one correction cancels the inverse error.
\writeup Factor $A+uv^*=A(I+A^{-1}uv^*)$ and apply Sylvester to obtain $c\det A$. Multiplying the proposed inverse on the left by $A+uv^*$ gives
\[I+uv^*A^{-1}-c^{-1}u(1+v^*A^{-1}u)v^*A^{-1}=I.
\]
A right inverse of a square matrix is its inverse, proving the formula. The known inverse update needs two matrix-vector products and an outer product, with $O(n^2)$ arithmetic operations instead of dense inversion's $O(n^3)$.
\end{solution}
\begin{solution}{linear:ex:hadamard4}
\idea Row orthogonality forces equal counts of all four sign pairs beneath a constant first row.
\writeup Column sign changes make row one all ones without changing orthogonality. Let $a,b,c,d$ count pairs $(+,+),(+,-),(-,+),(-,-)$ in rows two and three. Their pairings with row one give $a+b=c+d$ and $a+c=b+d$; pairing the two rows gives $a+d=b+c$. Together with the sum $n$, these equations force $a=b=c=d=n/4$. Thus $4\mid n$.
\end{solution}
\begin{solution}{linear:ex:maxdet3}
\idea Integer row operations expose divisibility that the real volume bound alone misses.
\writeup Add row one to every other row; all entries in those $n-1$ rows are even. Factoring two from each proves divisibility. At size three the determinant magnitude is a multiple of four below eight, since $3^{3/2}<8$. It is therefore at most four. The matrix $\begin{pmatrix}1&1&1\\1&-1&1\\1&1&-1\end{pmatrix}$ has determinant four by two row subtractions and triangular expansion, attaining the bound.
\end{solution}
\begin{solution}{linear:ex:minkowskieq}
\idea Equality in each normalised AM--GM step makes all relative eigenvalues equal.
\writeup After congruence set $A=I,B=\operatorname{diag}(d_i)$ with all $d_i>0$. The two AM--GM deficits are non-negative, so equality in their sum forces both to vanish. In particular $1/(1+d_i)$ is constant, hence $d_i=c>0$. Undoing congruence yields $B=cA$. Conversely determinant homogeneity gives equality for every such pair, including the automatic one-dimensional case.
\end{solution}
\begin{solution}{linear:ex:logdetderiv}
\idea A Hermitian congruence makes the perturbation eigenvalues real and the logarithm scalar.
\writeup Put $K=A^{-1/2}HA^{-1/2}$, with eigenvalues $\mu_i$. For small real $t$, $A+tH$ stays positive definite and $\log\det(A+tH)=\log\det A+\sum_i\log(1+t\mu_i)$. Its first derivative is $\sum\mu_i=\operatorname{tr}(A^{-1}H)$ and second is $-\sum\mu_i^2=-\operatorname{tr}((A^{-1}H)^2)\le0$. This applies at every point along a positive segment, proving concavity.
\end{solution}
\begin{solution}{linear:ex:schurmin}
\idea Removing the smallest eigenvalue leaves a positive matrix whose Schur product stays positive.
\writeup For $\lambda=\lambda_{\min}(A)\ge0$, Schur positivity gives $(A-\lambda I)\circ B\succeq0$. Hence $A\circ B\succeq\lambda\operatorname{diag}(b_{ii})\succeq\lambda\min_ib_{ii}\,I$. Eigenvalue monotonicity proves the bound.
\end{solution}
\begin{solution}{linear:ex:kernels}
\idea Each entry is an inner product or a positive-coefficient function of a Gram entry.
\writeup Indicators $\mathbf1_{[0,s_i]}$ in $L^2[0,\infty[$ have Gram entry $\min(s_i,s_j)$. Functions $e^{-x_ir}$ have Gram entry $\int_0^\infty e^{-(x_i+x_j)r}\dd r=(x_i+x_j)^{-1}$. Planar unit vectors $(\cos\theta_i,\sin\theta_i)$ give the cosine matrix. Finally the Gram entries $u_i\cdot u_j$ have absolute value below one, so the convergent entrywise geometric series and Schur positivity give $(1-u_i\cdot u_j)^{-1}$. All four arguments also test complex coefficients through the usual Gram quadratic form.
\end{solution}
\begin{solution}{linear:ex:gershscale}
\idea Keep the remote eigenvalue discs away while making the first radius quadratic in the perturbation.
\writeup The off-diagonal perturbation has eigenvalues $2\varepsilon,-\varepsilon,-\varepsilon$, so Weyl gives $\mu\le1+2\varepsilon<3/2$. For $D=\operatorname{diag}(1,4\varepsilon,4\varepsilon)$, the similar matrix has row radii $8\varepsilon^2,1/4+\varepsilon,1/4+\varepsilon$. The discs centred at two and three lie strictly to the right of $3/2$, because $\varepsilon<1/4$. They cannot contain $\mu$, so the disc at one must contain it, proving $|\mu-1|\le8\varepsilon^2$. This is an exact bound, with no numerical perturbation estimate required.
\end{solution}
\begin{solution}{linear:ex:diagdomdet}
\idea Along a strictly dominant deformation the real determinant cannot change sign without vanishing.
\writeup If $D$ is the diagonal part, $A_t=D+t(A-D)$ stays strictly dominant for $0\le t\le1$ and therefore invertible. Its determinant is a continuous real polynomial in $t$, nowhere zero. It keeps the positive initial sign $\det D=\prod_ia_{ii}>0$, yielding $\det A>0$.
\end{solution}
\begin{solution}{linear:ex:oppenheim}
\idea Compare the first Schur complement with a smaller Schur product, then use induction and determinant monotonicity.
\writeup Approximate $A$ by $A+\varepsilon I$ and pass to the limit at the end, so assume $A\succ0$. Size one is equality. Write $A=\begin{pmatrix}a&\alpha^*\\\alpha&A_1\end{pmatrix}$ and $B=\begin{pmatrix}b&\beta^*\\\beta&B_1\end{pmatrix}$. Then $S=A_1-\alpha\alpha^*/a\succ0$ and $\det A=a\det S$. If $b=0$, positivity forces $\beta=0$ and both sides vanish. Otherwise $B_1-\beta\beta^*/b\succeq0$, and the Schur complement in the product is
\[A_1\circ B_1-\frac{(\alpha\circ\beta)(\alpha\circ\beta)^*}{ab}
 =S\circ B_1+\frac1a(\alpha\alpha^*)\circ(B_1-\beta\beta^*/b)\succeq S\circ B_1.
\]
Both terms are positive by Schur's theorem. Determinant monotonicity and induction give $\det(A\circ B)\ge ab\det(S\circ B_1)\ge ab\det S\prod_{i\ge2}b_{ii}=\det A\prod_ib_{ii}$. Continuity proves the remaining semidefinite cases.
\end{solution}
\begin{solution}{ex:v8cEa}
\idea A nonzero Hermitian direction has a strictly negative second log-determinant derivative.
\writeup Under $H\mapsto A^{-1/2}HA^{-1/2}$, the form becomes $\operatorname{tr}(H'K')$, a real symmetric bilinear form on Hermitian matrices. At $H=K$ its value is $\operatorname{tr}((H')^2)=\sum_i\lambda_i(H')^2$, positive unless $H=0$. Thus it is a real inner product and the second derivative is $-g_A(H,H)$. Along a segment with $A\ne B$, its fixed nonzero direction has strictly negative second derivative everywhere. The scalar strict-concavity criterion gives strict inequality at each interior parameter.
\end{solution}
\begin{solution}{ex:v8cEb}
\idea A positive Fourier density makes a zero quadratic form force a vanishing finite exponential sum.
\writeup The Gaussian transform from Chapter~\ref{ch:fourier} gives
\[e^{-\gamma\|x-y\|^2}=\int_{\R^d}(\pi/\gamma)^{d/2}e^{-\pi^2\|\xi\|^2/\gamma}e^{2\pi i\xi\cdot(x-y)}\dd\xi.
\]
For coefficients $c_i$, their kernel quadratic form is the integral of the squared modulus of a finite exponential sum against this strictly positive density. If zero, continuity makes that sum zero everywhere. Choose $v$ with all $v\cdot x_i$ distinct: for $d\ge1$ the choice $v=(1,t,\ldots,t^{d-1})$ avoids a finite set of polynomial roots. Along $\xi=sv$, differentiating at $s=0$ through orders zero to $N-1$ gives a Vandermonde system in the distinct projected values. Thus every $c_i=0$. For a single point positivity is immediate, also covering dimension zero's only possible distinct-point case.
\end{solution}
\begin{solution}{ex:v8cEc}
\idea Reduce the determinant to the update dimension and cancel the inverse correction by the same small matrix.
\writeup Factor $A$ and apply Sylvester with $A^{-1}U,V^*$ to get the determinant identity. In the product of $A+UV^*$ with the proposed inverse, the correction terms combine as
\[UV^*A^{-1}-U(I_r+V^*A^{-1}U)C^{-1}V^*A^{-1}=0.
\]
The product is identity, proving the inverse formula. The determinant identity also shows that invertibility of the update is equivalent to invertibility of $C$.
\end{solution}
\begin{solution}{ex:v8cEd}
\idea A nearly scale-invariant integral supplies a Rayleigh quotient as close to $\pi$ as desired.
\writeup On $[m+1,m+2]\times[n+1,n+2]$, the positive integrand is at most $1/(\sqrt{(m+1)(n+1)}(m+n+2))$, hence at most the corresponding sum term. Summing unit-square bounds proves the first inequality. Write $L=N+1$. The restricted domain lies in $[1,L]^2$, and substitution $y=xt$ gives its integral exactly as
\[(\log L-2\log M)\int_{1/M}^M\frac{\dd t}{\sqrt t(1+t)}.
\]
The squared norm of the coefficient vector is $H_N\sim\log N$. Its Rayleigh quotient therefore has limit inferior at least this last $t$ integral for each fixed $M$. Letting $M\to\infty$, substitute $t=u^2$ to make the limit $\int_0^\infty2/(1+u^2)\dd u=\pi$. The proved upper bound is $\pi$, so $\lambda_{\max}(H_N)\to\pi$, establishing sharpness.
\end{solution}

\section*{Sources and further reading}

\begin{itemize}
\item R.~Egan, I.~Kotsireas, S.~London and P.~\'O Cath\'ain, \emph{Recent progress in the construction of Hadamard matrices}, arXiv:2610.27175, version 2, 3 November 2026.
\item J.~Hadamard, \emph{R\'esolution d'une question relative aux d\'eterminants}, Bulletin des Sciences Math\'ematiques 17 (1893),
240--246; H.~Kharaghani and B.~Tayfeh-Rezaie, \emph{A Hadamard matrix of order 428}, Journal of Combinatorial Designs 13 (2005), 435--440.
\item R.~Bhatia, \emph{Matrix Analysis}, Springer, 1997, Chapters V (operator monotone functions, L\"owner's theorem) and IX; R.~Bhatia,
\emph{Positive Definite Matrices}, Princeton University Press, 2007.
\item R.~A.~Horn and C.~R.~Johnson, \emph{Matrix Analysis}, 2nd ed., Cambridge University Press, 2013, Chapters 6 (Gershgorin) and 7
(positive definite matrices, Hadamard, Fischer, Oppenheim, Schur product).
\item I.~Schur, \emph{Bemerkungen zur Theorie der beschr\"ankten Bilinearformen mit unendlich vielen Ver\"anderlichen}, Journal f\"ur die
reine und angewandte Mathematik 140 (1911), 1--28 (Schur product theorem; Hilbert's inequality with the sharp constant $\pi$).
\item G.~H.~Hardy, J.~E.~Littlewood and G.~P\'olya, \emph{Inequalities}, Cambridge University Press, 1934, Chapter IX (Hilbert's inequality);
J.~M.~Steele, \emph{The Cauchy--Schwarz Master Class}, Cambridge University Press, 2004, Chapter 10.
\item S.~Gershgorin, \emph{\"Uber die Abgrenzung der Eigenwerte einer Matrix}, Izvestiya Akademii Nauk SSSR 6 (1931), 749--754; R.~S.~Varga,
\emph{Ger\v{s}gorin and His Circles}, Springer, 2004.
\item K.~L\"owner, \emph{\"Uber monotone Matrixfunktionen}, Mathematische Zeitschrift 38 (1934), 177--216; E.~Heinz, \emph{Beitr\"age zur
St\"orungstheorie der Spektralzerlegung}, Mathematische Annalen 123 (1951), 415--438.
\item B.~Sch\"olkopf and A.~J.~Smola, \emph{Learning with Kernels}, MIT Press, 2002, Chapter 2 (positive definite kernels).
\item T.~M.~Cover and J.~A.~Thomas, \emph{Elements of Information Theory}, 2nd ed., Wiley, 2006, Chapter 17 (determinant inequalities
from entropy).
\end{itemize}
